\chapter{Technical interfaces}
\label{pa-app-technical}

The main chapters state the etale Picard--Albanese route.  These interface
nodes retain the nontrivial formalization bridges while leaving routine
transport and standard library facts in the companion record.

\begin{proposition}[Cohomology interface]
  \label{pa-app-cohomology-interface}
  \dcref{III.1,III.12,custom}
  \uses{pa-thm-cohomology}
  \cite{stacks-project}
  Cech cohomology on affine covers computes the groups used for the tangent
  space of the Picard functor and commutes with field extension.
\end{proposition}

\begin{proof}
  Affine acyclicity reduces derived cohomology to a bounded Cech complex.
  Tensoring with a field is flat and commutes with its terms, so it commutes
  with the resulting kernels and cokernels.
\end{proof}

\begin{proposition}[Etale Picard interface]
  \label{pa-app-etale-picard-interface}
  \dcref{2.2,4.8,custom}
  \uses{pa-thm-picard-representability,pa-prop-picard-group}
  \cite{kleiman-picard}
  The relative Picard presheaf, after etale sheafification, is represented by
  the quotient obtained from the bounded divisor parameter space.
\end{proposition}

\begin{proof}
  Rigidified line bundles remove pullbacks from the base.  The invertible locus
  in the bounded parameter space is open, and effective descent identifies its
  sheafification with the relative Picard functor.
\end{proof}

\begin{proposition}[Divisor and symmetric-power interface]
  \label{pa-app-divisor-interface}
  \dcref{3.7,III.6,custom}
  \uses{pa-lem-symmetric-power-represents,pa-def-symmetric-power,pa-def-abel-jacobi}
  \cite{kleiman-picard,milne-av}
  Finite flat effective divisors on a smooth curve are represented by its
  symmetric powers, and the universal divisor defines the Abel--Jacobi map.
\end{proposition}

\begin{proof}
  A finite flat divisor is locally cut out by a monic polynomial; its roots
  give the symmetric-power coordinates.  Tensoring the associated divisor
  line bundles is invariant under permutation and descends to the quotient.
\end{proof}

\begin{proposition}[Picard identity-component interface]
  \label{pa-app-pic0-interface}
  \dcref{5.1,5.4,5.11,custom}
  \uses{pa-thm-pic0-abelian,pa-prop-picard-tangent}
  \cite{kleiman-picard}
  The identity component of the represented Picard scheme is an abelian
  variety, and its tangent dimension is \(\dim_k H^1(C,\mathcal O_C)=g(C)\).
\end{proposition}

\begin{proof}
  Degree is locally constant, so the identity component is open and closed.
  The tangent calculation gives the dimension, while the large-degree divisor
  construction gives properness; translation supplies smoothness everywhere.
\end{proof}

\begin{theorem}[Albanese interface]
  \label{pa-app-albanese-interface}
  \dcref{I.2,III.6,3.1,3.2,custom}
  \uses{pa-thm-rigidity,pa-thm-rational-extension,pa-thm-albanese}
  \cite{milne-av,mumford-av}
  Rigidity and codimension-one extension give the universal factorisation of a
  pointed curve map through the Albanese of the curve.
\end{theorem}

\begin{proof}
  Symmetric sums of a pointed map descend from ordered divisors.  Rigidity
  makes the resulting map constant on complete linear-system fibres, and the
  extension theorem removes the rational-map indeterminacy.
\end{proof}

\begin{theorem}[Jacobian interface]
  \label{pa-app-jacobian-interface}
  \dcref{III.6,5.13,5.14,custom}
  \uses{pa-thm-jacobian,pa-thm-jacobian-base-change}
  \cite{kleiman-picard}
  The Jacobian is the degree-zero etale Picard component, with its group
  structure, universal map, and compatibility with base change.
\end{theorem}

\begin{proof}
  The etale Picard representation restricts to the open-and-closed identity
  component.  The universal line bundle gives the Abel--Jacobi map, and the
  same construction after field extension identifies the two Jacobians.
\end{proof}
